\section{Introduction}
\label{sec:intro}

\begin{figure}[h!]
  \centering
  \includegraphics[width=1.0\linewidth]{figs/0_teaser.pdf}
  \vspace{-6mm}
  \caption{\textbf{Human-centered scene reconstruction and free-view video synthesis.} (a) Source view inputs, (b) our metric scale point map reconstruction, and (c) free-view rendering with our feed-forward Gaussian Splatting.}
  \vspace{-6mm}
  \label{fig:teaser}
\end{figure}

Feed-forward free-viewpoint video synthesis is a crucial task, especially in the setting of sparse views, which could serve many downstream applications such as telecommunications, stage/sports broadcasts, and so on.
%
Existing pipelines are typically based on differentiable rendering~\cite{mildenhall2020nerf, xu2022point-nerf, yu2021pixelnerf, chen2021mvsnerf, lin2022enerf, Wu2024_4dgs, sun20243dgstream} with the development of neural network. 
%
Particularly, Gaussian Splatting~\cite{kerbl2023_3dgs} shows its advancement for the high efficiency of rendering and the capable mechanism of back-propagation, but relies on minute-level optimization for each scene and very dense input views.


Recently, Gaussian related methods~\cite{chen2024mvsplat, liu2024mvsgaussian, charatan2023pixelsplat} achieve instant inference in a feed-forward manner, avoiding per-scene optimization, for real-time applications such as telecommunication systems~\cite{tu2024tele} and human-scene synthesis~\cite{zheng2024gpsgaussian, zhou2024gps}.
%
However, a common strategy is to estimate Gaussian primitive maps defined on the source views, leveraging geometry proxies of multi-view stereo (MVS)~\cite{chen2024mvsplat, liu2024mvsgaussian} or binocular stereo-matching~\cite{zheng2024gpsgaussian, zhou2024gps}.
%
Such methods require a large overlap of paired images, which increases the redundancy of Gaussian primitives in this overlap area. 
%
Otherwise, they could not provide reasonable geometry prior when input cameras are with large sparsity.


More recently, DUSt3R~\cite{wang2024dust3r, leroy2024grounding} proposes a novel geometry representation as point maps of input views, which assigns each pixel of input images to a free 3D point.
%
Unlike traditional MVS methods~\cite{yao2018mvsnet, yang2020cost}, DUSt3R gets rid of stereo constraints and achieves pixel-aligned point maps of binocular inputs from very sparse views, by training on immense 3D geometry data.
%
To alleviate the impact of the high degree of freedom, DUSt3R and its follow-ups~\cite{leroy2024grounding, wang2024moge} normalize the scale of the reconstructed point map with an average point distance of each scene, which typically causes a dramatic instability of reconstruction in consecutive frames, see Fig.~\ref{fig:teaser}(b). 
%
Some recent methods~\cite{smart2024splatt3r, ye2024nopose} leverage point map representation for static scene rendering in canonical space.
However, human movement in the scene would cause a relative depth difference in canonical space and lead to large jitters for free-view video rendering, due to the lack of stereo constraint.
%
In addition, training a foundation model for scale-aware geometry typically requires immense 3D data, but it is always tedious and cumbersome to acquire 3D geometry data.
%
Therefore, the key point is to obtain scale-aware geometry in a self-supervised manner and to ease the burden of 3D data acquisition.


In this paper, we propose Splat-SAP to achieve human-centered scene reconstruction in real metric space and feed-forward rendering of novel views via Gaussian plane, when inputting a pair of images and camera calibration.
%
Unlike \cite{wang2024dust3r, leroy2024grounding} representing a scale-invariant point map in canonical space, we inject camera intrinsic embedding~\cite{ye2024nopose} and global image feature~\cite{wang2025vggt} into a network as input to learn a \textbf{scaling} factor to transform the estimated point map from canonical space to real space.
%
Since the point representation in the original design of DUSt3R~\cite{wang2024dust3r, leroy2024grounding} lacks stereo constraints between two source views, there always exists misalignment between the two point maps.
%
Thus, we compute the cost between 2 source views to do an iterative \textbf{\textit{coarse registration}} of 2 reconstructed point maps, by projecting the feature from one view to another with calibrated camera pose.
%
This registration is in the format of a translation map, denoting pixel-wise shift.
%
Our scaling factor by intrinsic embedding and translation learning by extrinsic projection compose exactly an \textbf{affine transformation} of point map from canonical space to real space.


In terms of rendering, we anchor Gaussian primitives directly on the target view as a Gaussian plane, so that reducing the redundancy of using directly two point maps of source views as Gaussian positions~\cite{zheng2024gpsgaussian, chen2024mvsplat, smart2024splatt3r, ye2024nopose}.
%
Depth of the Gaussian plane is initialized by projecting two point maps via $\alpha$-blending~\cite{kerbl2023_3dgs}, which largely eases the burden of accurate depth estimation.
%
Further, we do a \textbf{\textit{fine registration}} with strict stereo constraint~\cite{yao2018mvsnet, lin2022enerf, liu2024mvsgaussian}, which relies on a 3D cost volume by sampling several depth candidates along each pixel ray near the initialized depth. 
%
With such 3D cost representation, we can aggregate more 3D information to overcome the unobservation issue due to the large sparsity, and to estimate accurate depth.  
%
The color of the Gaussian plane can be initialized by warping source view pixels directly via the estimated depth.
%
Additionally, we incorporate both fine 2D and dense 3D features to estimate Gaussian primitives and refine Gaussian color for high-quality rendering.


More importantly, our pipeline can be trained without geometry supervision, different from \cite{wang2024dust3r, leroy2024grounding, wang2024moge}.
%
To this end, we collect large-scale multi-view data of over 10,000 frames of motion sequences of human-centered scenes to train our model.
%
We validate the effectiveness of our method on diverse camera settings, \textit{e.g.} industry camera, mobile phone, and GoPro, for both reconstruction and rendering tasks.
%
In summary, we claim three following contributions:
\begin{itemize}
    \item We introduce a feed-forward pipeline to reconstruct scale-aware point maps and to render free-view video of human-centered scenes, where the point maps are trained in a self-supervised manner without any 3D supervision. 
    \item We propose a 2D coarse to 3D fine registration strategy to estimate scale-aware point maps with a learnable affinity.
    \item We design a Gaussian plane, leveraging scale-aware point maps and incorporating both 2D and 3D features, to guarantee the efficiency and completeness of rendering.
\end{itemize}
